\documentclass[11pt,a4paper]{article}
\usepackage[margin=27mm]{geometry}
\usepackage{amsmath,amssymb,amsthm,hyperref}
\newtheorem{lemma}{Lemma}
\newtheorem{theorem}[lemma]{Theorem}
\newcommand{\Sz}{\operatorname{Sz}}
\newcommand{\F}{\mathbb F}
\newcommand{\gen}[1]{\langle#1\rangle}
\title{Scalar fields and split torus families}
\author{Human Compression project}\date{}
\begin{document}\maketitle
\begin{abstract}
One elementary divisibility identity determines all intersections in a
family of Suzuki coset pregeometries. The same identity proves that the
cyclic scalar group in each affine residue acts irreducibly. This gives
a uniform proof of the intended families in Constructions 5 and 7 of the
fixed higher-rank source, after stating its necessary indexing and order
corrections explicitly. Only residual weak primitivity is asserted here.
\end{abstract}

\section{The arithmetic of the scalar subgroup}
Let $e\ge3$ be odd, let $k\ge1$, let $a_1,\ldots,a_k$ be distinct odd primes, and put
\[
 t=2^e,\quad A=\prod_{i=1}^k a_i,\quad q=t^A.
\]
Choose distinct primes $\ell_i$ such that
\[
 \ell_i\mid\frac{t^{a_i}-1}{t-1},\qquad
 M=(t-1)\prod_{i=1}^k\ell_i.
\]
The primes $\ell_i$ need not be coprime to $t-1$.
For $J\subseteq\{1,\ldots,k\}$ set
\[
 D_J=\prod_{i\notin J}a_i,\quad q_J=t^{D_J},\quad
 M_J=(t-1)\prod_{i\notin J}\ell_i.
\]

\begin{lemma}\label{arithmetic}
For every $J$ one has $\gcd(M,q_J-1)=M_J$. In particular $M\mid q-1$.
The multiplicative subgroup of order $M_J$ in $\F_{q_J}^*$ generates
$\F_{q_J}$ as a field over $\F_2$.
\end{lemma}
\begin{proof}
If $\ell_i\nmid t-1$, the order of $t$ modulo $\ell_i$ is $a_i$:
it divides that prime and is not one. Hence
$\ell_i\mid t^{D_J}-1$ exactly when $i\notin J$. Its exponent in $M$
is one, since the chosen $\ell_i$ are distinct.

If $\ell_i\mid t-1$, reduction of
$1+t+\cdots+t^{a_i-1}$ modulo $\ell_i$ shows that $\ell_i=a_i$.
For an odd prime $\ell\mid t-1$, successive binomial expansions give
\[
 v_\ell(t^D-1)=v_\ell(t-1)+v_\ell(D).
\]
For completeness, raising a number congruent to one modulo $\ell$ to
a prime power step $r$ preserves the valuation when $r\ne\ell$:
the first binomial term has coefficient $r$, and all later terms have
larger valuation. When $r=\ell$, the first term gains exactly one
power of $\ell$ and all later terms have strictly larger valuation,
since $\ell$ is odd. Iterating proves the displayed formula.
Since $D_J$ is squarefree, the additional factor $\ell_i$ in $M$ occurs
in $q_J-1$ precisely when $a_i\mid D_J$. All other prime powers of
$t-1$ occur on both sides. This proves the gcd identity.

Suppose that the subgroup of order $M_J$ is contained in a subfield
$\F_{2^f}\subseteq\F_{q_J}$. Its subgroup of order $2^e-1$ gives
$2^e-1\mid2^f-1$, so $e\mid f$ by the usual gcd identity for powers
of two. Write $f=eD$, where $D\mid D_J$. Applying the identity just
proved to $t^D-1$ shows that each $a_i$ with $i\notin J$ divides $D$;
otherwise its additional prime factor in $M_J$ would be missing.
Thus $D=D_J$, and the subfield is the whole field.
\end{proof}

\section{All subgroup intersections and the maximality condition}
We use the classical Suzuki matrix model, its standard subfield groups,
and maximality of prime-relative-degree inclusions of simple Suzuki
groups as attributed inputs. Write $S_d=\Sz(2^d)$ for a standard
subfield group, $Z_{2^d}$ for the involution centre of its standard
Sylow two-subgroup, and $H_r$ for the subgroup of order $r$ in the
ambient split torus. The standard Weyl involution $w$ inverts that
torus. Unique matrix coordinates give
\[
 S_b\cap S_c=S_{\gcd(b,c)}.
\]
Conjugation by the split torus on $Z_q$ is multiplication by a power
coprime to $q-1$. These coordinate and action facts have complete
derivations, with the classical ambient model explicit, in the acquired
Suzuki preliminary and normalizer notes. We do not assume the source's
classification of all residually weakly primitive pregeometries.

In $G=S_{eA}$ consider the following $k+2$ parabolics, with separate
labels for the last two:
\[
 P_i=S_{eA/a_i}\ (1\le i\le k),\qquad
 P_D=H_M\rtimes\gen{w},\qquad P_E=Z_q\rtimes H_M.
\]
For a set of selected subfield labels $J$, their intersection is
$S_{eD_J}$. Intersecting further with either or both of the last two
parabolics gives exactly
\begin{align*}
 S_{eD_J}\cap P_D&=H_{M_J}\rtimes\gen{w},\\
 S_{eD_J}\cap P_E&=Z_{q_J}\rtimes H_{M_J},\\
 S_{eD_J}\cap P_D\cap P_E&=H_{M_J}.
\end{align*}
Indeed the diagonal parameter must belong to $\F_{q_J}$, so Lemma
\ref{arithmetic} gives its precise cyclic order. The central unipotent
parameter is restricted to that same field. The reflection coset of
the split normalizer is outside the point stabilizer, leaving only
the torus in the third intersection. This also proves the formulas
when $J$ is empty.

\begin{theorem}\label{family}
The displayed parabolics define a residually weakly primitive coset
pregeometry of rank $k+2$. Its full intersection is $H_{t-1}$.
\end{theorem}
\begin{proof}
The condition is that every proper subset of the labels admits a
remaining parabolic whose intersection is maximal proper in the
current intersection. We check all four possibilities concerning
the labels $D,E$.

If neither has been selected and a subfield label remains, adjoining
any such label gives a prime-relative-degree inclusion
$S_{eD_J/a_i}<S_{eD_J}$. Both are simple, since their exponents are
at least $e\ge3$, so this is maximal. If no subfield label remains,
the current intersection is $S_e$, and adjoining $D$ gives its full
split torus normalizer, a classical maximal subgroup.

If $D$ but not $E$ has been selected, adjoining $E$ passes from
$H_{M_J}\rtimes\gen{w}$ to its cyclic subgroup of index two.

If $E$ but not $D$ has been selected, adjoining $D$ passes from
$Z_{q_J}\rtimes H_{M_J}$ to $H_{M_J}$. This is maximal because
the scalar action is irreducible over $\F_2$. To verify this last
assertion, an additive subspace invariant under the cyclic scalar
group is invariant under the field it generates. Lemma
\ref{arithmetic} identifies that field with $\F_{q_J}$, whose
one-dimensional vector space has only the zero and full subspaces.
Every intermediate subgroup containing $H_{M_J}$ is the product of
that complement and its intersection with $Z_{q_J}$; hence no proper
nontrivial intermediate subgroup exists.

Finally, if both $D,E$ have been selected, a subfield label $i$ remains
because the chosen subset is proper. Adjoining it passes from
$H_{M_J}$ to $H_{M_J/\ell_i}$, of prime index $\ell_i$.
This exhausts the possibilities and proves the condition. The
intersection formulas give the stated full intersection.
\end{proof}

\section{The two source families and their literal corrections}
The source is D.~Leemans, \emph{The residually weakly primitive
pre-geometries of the Suzuki simple groups}, Note di Matematica
20(1) (2000/2001), 1--20. We use its fixed twenty-page edition.
The introduction to the constructions, on p.~15, instructs the reader
to retain the intended whole intersection lattice from the rank-four
discussion. We therefore distinguish this intended lattice from
literal indexing or order errors in the displayed recipes.

For Construction 7 put $k=\alpha-2$, and take its selected relative
prime degrees as $a_1,\ldots,a_k$. Its definition of $t$ gives our
$e=\log_2t$. Its range $\alpha\le\sum_j e_j$ leaves at least two
prime factors, with multiplicity, in $e$, so $e\ge3$. Its cyclic
parameters are our $\ell_i$. The two displayed soluble group orders
are exactly those in Theorem~\ref{family}. The source labels the
selected Suzuki groups $G_1,\ldots,G_{\alpha-2}$ and then also labels
the dihedral group $G_{\alpha-2}$, while its final tuple begins with
$G_0$. Relabel the subfield groups $P_1,\ldots,P_k$ and keep $D,E$
as separate labels. With that explicit indexing correction, every
permitted parameter is realized, with all intended intersections and
the maximality condition proved above.

Construction 5 has the same family with $e$ prime. In its first case,
the source exponent is a product of distinct primes. The distinguished
prime is $e=p_i$, and the other $m-1$ primes are our selected $a_j$;
thus $k=m-1$ and the rank is $m+1$. In its second case, one prime
$p_m$ is repeated in the ambient exponent. Take $e=p_m$ and select
all $m$ distinct primes as the $a_j$, including $p_m$ itself. They
remain pairwise distinct; no coprimality between $e$ and $A$ is needed
anywhere in the proof. The rank is $m+2$.

The literal recipes in Construction 5 require more than relabeling.
Both displayed dihedral orders omit the factor $t-1$. Its first case
assigns conflicting indices. Its second case uses an undefined $i$
in the affine product and ends its tuple at $G_m$ despite listing
$G_{m+1}$. The corrected dihedral and affine groups are, in both
cases, $D_{2M}$ and $Z_q\rtimes H_M$, with the product taken over
all selected indices. These are the groups required by the stated
rank-four intersections in the preceding discussion. We prove this
corrected family, not the inconsistent literal assignment.

The method extends the previously retained rank-three split family:
subfield coordinates give the intersections, and scalar irreducibility
proves the affine maximality step. The uniform arithmetic lemma supplies
the higher-rank argument, including the case $\ell_i\mid t-1$.
No priority claim is made for this elementary arithmetic or for the
classical Suzuki subgroup structure. Arbitrary relative positions,
flag transitivity, residual connectedness, a classification of all
families, and compression measurements are outside this component.
\end{document}
